% !TeX root = ../../Book5.tex
% 纲要标题：### 8.6 Brauer 对应与 Green 对应
% 纲要位置：工作记录/计划纲要.md（2026-10-20 新增骨架）。
% * 状态：初稿已完成；新增范围尚未正式审核。
% 写作范围：依次标明 Brauer 主定理及 Green 对应的假设、陈述、证明概要与引用范围，不承担深层块分类及研究猜想。
% #### 8.6.1 缺陷群的共轭性与 Brauer 对应
% #### 8.6.2 Brauer 三个主定理的陈述与证明范围
% #### 8.6.3 Green 对应
% #### 8.6.4 小群的块与对应计算

\section{Brauer 对应与 Green 对应}
\label{b5:sec:0806}

对一个块取 Brauer 像，得到的是较小群代数中的幂等元；把一个不可分解模限制到子群，得到的模却可能分裂。这两个操作都有选择问题：哪个局部块应当与原块对应，限制模中的哪一项应当与原模对应？缺陷群为第一个问题指定子群，顶点为第二个问题指定直和项，后者由 Green 对应实现\footnote{格林（James Alexander Green，1926-2014），英国数学家，研究有限群表示，建立不可分解模的 Green 对应及有限一般线性群的特征标理论。}。本节分别说明这两个选择怎样成立，并用 \(S_3\) 比较两种对应。涉及块的主定理时，假设特征 \(p\) 的域 \(k\) 对所讨论群及其全部子群都是分裂域；取 \(k\) 代数闭可满足这一要求。单纯关于顶点和 Green 对应的结论不需要分裂假设。

\subsection{缺陷群的共轭性与 Brauer 对应}
\label{b5:subsec:0806:01}

\cref{b5:def:block-defect-group}用 \(\operatorname{Br}_D^G(b)\ne0\) 的极大性定义缺陷群。它便于从块幂等元的系数计算，却没有立即说明两个极大子群为什么共轭。把块同时当作左模和右模，就可以利用顶点的唯一性。具体地，令 \(G\times G\) 在 \(kG\) 上作用为
\[
	(g,h)\cdot a=gah^{-1},\qquad
	\Delta H=\{\,(h,h)\mid h\in H\,\}\le G\times G.
\]
\(G\times G\) 的子模恰是双侧理想；块 \(kGb\) 因而是不可分解的 \(k[G\times G]\)-模。这里一定要保留两侧作用：块作为左 \(kG\)-模本来就是投射模，其左模顶点为一，不能记录缺陷群。

\begin{theorem}\label{b5:thm:defect-vertices-traces}
	设 \(G\) 为有限群，\(k\) 为特征 \(p>0\) 的域，\(b\) 为 \(kG\) 的块幂等元。存在 \(p\)-子群 \(D\le G\)，使 \(\Delta D\) 为上述双模 \(kGb\) 的顶点。对任意这样的 \(D\)，有
	\[
		b=\operatorname{Tr}_D^G(a)\quad\text{对某个}\;a\in(kGb)^D,
	\]
	且对每个 \(p\)-子群 \(Q\le G\)，
	\[
		\operatorname{Br}_Q^G(b)\ne0
		\quad\Longleftrightarrow\quad
		Q\le {}^gD\quad\text{对某个}\;g\in G.
	\]
	特别地，\(b\) 的缺陷群恰为 \(D\) 的 \(G\)-共轭子群。
\end{theorem}
\begin{proof}
	映射 \((g,h)\otimes1\mapsto gh^{-1}\) 给出
	\[
		\operatorname{Ind}_{\Delta G}^{G\times G}k\cong kG.
	\]
	因此 \(kGb\) 相对 \(\Delta G\) 投射。由\cref{b5:thm:vertices-sources}，可以在 \(\Delta G\) 中选择它的一个顶点，必为 \(\Delta D\)。若 \(\Delta D_1=(g,h)\Delta D(g,h)^{-1}\)，投影到第一坐标就有 \(D_1=gDg^{-1}\)，所以这些 \(D\) 互相共轭。

	对任意 \(H\le G\)，先证明相对迹判据
	\[
		kGb\text{ 相对 }\Delta H\text{ 投射}
		\quad\Longleftrightarrow\quad
		b=\operatorname{Tr}_H^G(a)\text{，其中 }a\in(kGb)^H.
	\]
	正向先限制到 \(\Delta G\)。Mackey 分解中出现的交集为 \(\Delta G\cap(g,h)\Delta H(g,h)^{-1}\)；将 \((g,h)\) 写成 \((g,g)(1,g^{-1}h)\)，可见每个交集都包含于某个 \(\Delta G\)-共轭的 \(\Delta H\)。所以该限制模相对 \(\Delta H\) 投射。Higman 判据给出 \(1_{kGb}=\operatorname{Tr}_H^G(\alpha)\)，其中 \(\alpha\) 对 \(H\) 的共轭作用线性。将等式作用到被 \(G\) 固定的 \(b\)，便得 \(b=\operatorname{Tr}_H^G(\alpha(b))\)。反向若 \(b=\operatorname{Tr}_H^G(a)\)，则左乘 \(a\) 是共轭 \(H\) 模自同态，其相对迹为左乘 \(b\)，也就是恒等。因此限制到 \(\Delta G\) 的模相对 \(\Delta H\) 投射。原模已相对 \(\Delta G\) 投射，再利用诱导的传递性，即得原模相对 \(\Delta H\) 投射。

	取顶点 \(\Delta D\)，便有所需的迹表达式，且 \(D\) 对此性质极小。若 \(\operatorname{Br}_D^G(b)=0\)，由\subsecref{b5:subsec:0804:02}中按共轭轨道求得的核，可写
	\[
		b=\sum_{Q<D}\operatorname{Tr}_Q^D(u_Q),\qquad u_Q\in(kG)^Q.
	\]
	与 \(b=\operatorname{Tr}_D^G(a)\) 相乘，利用 \(a\) 的 \(D\)-不变性和迹的传递性，再乘一次中心元 \(b\)，得到
	\[
		b=\sum_{Q<D}\operatorname{Tr}_Q^G(ba u_Q).
	\]
	各项都在交换局部代数 \(bZ(kG)\) 中。这里 \(bZ(kG)\) 就是不可分解双模 \(kGb\) 的自同态环：双模同态由 \(b\) 的像决定，该像须与整个群代数交换；局部性因此由\cref{b5:lem:relative-intersection}证明中回顾的局部自同态环结论给出。一个局部环中若若干元素之和为单位，至少一项是单位。于是某个 \(\operatorname{Tr}_Q^G(ba u_Q)\) 在 \(bZ(kG)\) 中可逆，乘其中心逆元就把 \(b\) 写成从真子群 \(Q<D\) 来的相对迹，矛盾。因此 \(\operatorname{Br}_D^G(b)\ne0\)。

	对任意 \(p\)-子群 \(Q\)，将 \(\operatorname{Tr}_D^G(a)\) 按 \(Q\backslash G/D\) 分组，得到
	\[
		b=\sum_{x\in Q\backslash G/D}
		\operatorname{Tr}_{Q\cap{}^xD}^{Q}(xax^{-1}).
	\]
	若所有交集都是 \(Q\) 的真子群，Brauer 映射杀掉右侧每一项。因此像非零迫使某个 \(Q\le{}^xD\)。反过来，若 \(Q\le{}^xD\)，则 \(C_G({}^xD)\subseteq C_G(Q)\)；\(b\) 在前者上已有非零系数，在后者上的截断也非零。于是极大的 \(Q\) 正是 \(D\) 的各个共轭。
\end{proof}

这个证明还使缺陷群与块中的各个模发生联系。
\begin{corollary}\label{b5:cor:block-modules-relative-defect}
	设 \(G\) 为有限群，\(k\) 为特征 \(p>0\) 的域，块 \(kGb\) 的缺陷群为 \(D\)。若有限维左 \(kG\)-模 \(M\) 满足 \(bM=M\)，则 \(M\) 相对 \(D\) 投射。特别地，不可分解的 \(M\) 有一个顶点包含于 \(D\)。
\end{corollary}
\begin{proof}
	取 \(b=\operatorname{Tr}_D^G(a)\)。元素 \(a\) 在 \(M\) 上的作用是 \(kD\)-自同态，而它的相对迹是 \(b\) 的作用，即 \(1_M\)。由\cref{b5:thm:higman-criterion}得到相对投射性，再用顶点的包含性质即得后一句。
\end{proof}

现在固定 \(p\)-子群 \(D\)，令 \(N=N_G(D)\)。对 \(z\in Z(kG)\)，\(\operatorname{Br}_D^G(z)\) 不仅在 \(kC_G(D)\) 中，而且被 \(N\) 共轭固定，因此也属于 \(Z(kN)\)。若 \(1=\sum_i b_i\) 是 \(kG\) 的块分解，那么 \(1=\sum_i\operatorname{Br}_D^G(b_i)\) 是 \(Z(kN)\) 中的正交幂等分解，部分项可以为零。
\begin{definition}\label{b5:def:brauer-correspondent}
	设 \(G\) 为有限群，\(k\) 为特征 \(p>0\) 的域，\(D\le G\) 为 \(p\)-子群，\(N=N_G(D)\)，\(c\) 为 \(kN\) 的块幂等元。在 \(kG\) 的块幂等元中，恰有一个 \(b\) 满足
	\[
		c\operatorname{Br}_D^G(b)=c.
	\]
	记此 \(b\) 为 \(c^G\)，称为 \(c\) 的 \term[Brauer-duiying-kuai]{Brauer 对应块}[Brauer correspondent]。这里固定使用子群 \(N_G(D)\) 及映射 \(\operatorname{Br}_D^G\)。
\end{definition}
\symidx{\(c^G\)：子群块的 Brauer 对应块}
存在与唯一性只用 \(c\) 的本原性：\(c=\sum_i c\operatorname{Br}_D^G(b_i)\) 是正交中心幂等分解，恰有一项非零。此时尚不能断言 \(c\mapsto c^G\) 为双射，也不能断言两块具有相同缺陷群；第一主定理需要明确加入缺陷群恰为 \(D\) 的条件。

\subsection{Brauer 三个主定理的陈述与证明范围}
\label{b5:subsec:0806:02}

\term[Brauer-diyi-zhudingli]{Brauer 第一主定理}[Brauer's first main theorem]把缺陷群恰为 \(D\) 的块与其正规化子中的相应块配对。这里固定实际子群 \(D\)，块的缺陷群列表则仍按共轭理解。

\begin{theorem}[Brauer 第一主定理]\label{b5:thm:brauer-first-main}
	设 \(G\) 为有限群，\(k\) 为特征 \(p>0\) 的域，并且对 \(G\) 及其全部子群都是分裂域。给定 \(p\)-子群 \(D\le G\)，记 \(N=N_G(D)\)。则
	\[
		b\longmapsto\operatorname{Br}_D^G(b)
	\]
	给出缺陷群为 \(D\) 的 \(kG\) 块与缺陷群为 \(D\) 的 \(kN\) 块之间的双射，逆映射为 \(c\mapsto c^G\)。特别地，在这一范围内，\(\operatorname{Br}_D^G(b)\) 本身就是一个块幂等元。
\end{theorem}
\begin{proofsketch}
	相对迹与截断的相容公式是
	\[
		\operatorname{Br}_D^G\operatorname{Tr}_D^G(a)
		=\operatorname{Tr}_D^N\operatorname{Br}_D^G(a)
		\qquad(a\in(kG)^D).
	\]
	按 \(D\backslash G/D\) 分组后，\(D\cap{}^xD<D\) 的项都被截断杀掉；剩下的代表元恰来自 \(N/D\)。若 \(c\) 是 \(kN\) 中缺陷为 \(D\) 的块，令 \(b=c^G\)。利用这一公式和\cref{b5:thm:defect-vertices-traces}，可以把 \(c\) 放入
	\[
		\operatorname{Br}_D^G\bigl(b\operatorname{Tr}_D^G((kG)^D)\bigr)
	\]
	中。括号内是 \(bZ(kG)\) 的理想。它的像含非零幂等元，故该理想不可能是幂零理想；有限维局部代数的真理想全在幂零根中，所以它必须等于 \(bZ(kG)\)。于是 \(b\) 来自 \(D\) 的相对迹，而 \(\operatorname{Br}_D^G(b)\ne0\)，从而 \(D\) 正是它的缺陷群。同一个局部代数的像至多有一个非零幂等元，进一步给出 \(\operatorname{Br}_D^G(b)=c\)。

	这段论证还使用两个块论结论：正规 \(p\)-子群包含在每个块的缺陷群中，且块幂等元的支集包含在该正规子群的中心化子中。前者保证 \(kN\) 的块的缺陷群包含 \(D\)，后者允许在相对迹中先做截断。它们也使反向构造穷尽 \(kG\) 中缺陷为 \(D\) 的块。这里没有证明这两个一般结论；其证明分别见\cite[Corollary 12.3.4 and Proposition 12.5.7]{Webb2016FiniteGroupRepresentations}，上述局部理想论证及双射的完整证明见该书 Theorem 12.6.4。
\end{proofsketch}

\term[Brauer-dier-zhudingli]{Brauer 第二主定理}[Brauer's second main theorem]联系的是普通特征标在 \(p\)-奇异元素处的值与中心化子的 Brauer 特征标。设 \((K,\mathcal O,k)\) 是对 \(G\) 及其子群分裂的 \(p\)-模系统，\(\mathcal O\) 完备，且 \(K\) 含有 \(|G|\) 次单位根。给定 \(p\)-元素 \(u\in G\)，令 \(C=C_G(u)\)。对普通不可约特征标 \(\chi\)，函数 \(v\mapsto\chi(uv)\) 是 \(C\) 的 \(p\)-正则元素上的类函数，因而可唯一写成
\begin{equation}\label{b5:eq:generalized-decomposition-numbers}
	\chi(uv)=\sum_{\varphi\in\operatorname{IBr}(C)}d^u_{\chi\varphi}\varphi(v),
	\qquad v\in C_{p'}.
\end{equation}
其中使用\cref{b5:thm:brauer-character-basis}。普通特征标、Brauer 特征标及单位根都沿用\secref{b5:sec:0803}中固定的分圆域嵌入，统一取复数值。系数称为\term[guangyi-fenjie-shu]{广义分解数}[generalized decomposition number]。它们一般不是非负整数：将表示限制到 \(C\) 后，中心元 \(u\) 的各特征空间均是 \(KC\)-子模；每个子模的分解数为整数，再乘 \(u\) 的相应特征值并相加，就得到 \(d^u_{\chi\varphi}\)。因此这些系数是 \(p\) 次幂单位根的整数线性组合；\(u=1\) 时恢复普通分解数。
\symidx{\(d^u_{\chi\varphi}\)：广义分解数}

还须说明普通特征标“属于一个模特征块”的含义。由于 \(\mathcal O\) 完备，\(kG\) 的每个中心幂等元 \(b\) 唯一提升为 \(Z(\mathcal OG)\) 的幂等元 \(\widehat b\)。这里的提升可直接在中心中进行：共轭类和给出中心的基，因此约化映射 \(Z(\mathcal OG)\to Z(kG)\) 满射；从任一近似提升 \(a\) 出发，因 \((2a-1)^2\equiv1\pmod\pi\)，\(2a-1\) 在完备代数中可逆；以 \(a-(a^2-a)(2a-1)^{-1}\) 反复修正，误差的 \(\pi\)-进阶数加倍，完备性给出极限。两个同余的中心幂等提升之差为零，因为它们的交叉项 \(e(1-e')\)、\(e'(1-e)\) 都是落在 \(\pi\mathcal OG\) 中的幂等元，必须落在每个 \(\pi^n\mathcal OG\) 中，因而为零。各块的提升仍正交且和为一，所以一个不可约 \(KG\)-模上恰有一个 \(\widehat b\) 作用为恒等；它的普通特征标便称为属于块 \(b\)。

\begin{theorem}[Brauer 第二主定理]\label{b5:thm:brauer-second-main}
	设 \(G\) 为有限群，\((K,\mathcal O,k)\) 为对 \(G\) 及其全部子群分裂的 \(p\)-模系统，其中 \(\mathcal O\) 完备，且 \(K\) 含有 \(|G|\) 次单位根。取 \(p\)-元素 \(u\in G\)、普通不可约特征标 \(\chi\)，以及 \(C=C_G(u)\) 的不可约 Brauer 特征标 \(\varphi\)。若 \(b\) 为 \(\chi\) 所属块在 \(kG\) 中的幂等元，\(c\) 为 \(\varphi\) 所属 \(kC\) 块的幂等元，则
	\[
		d^u_{\chi\varphi}\ne0
		\quad\Longrightarrow\quad
		c\operatorname{Br}_{\langle u\rangle}^G(b)=c.
	\]
	也就是说，\eqref{b5:eq:generalized-decomposition-numbers}只有位于该 Brauer 像之下的局部块可能贡献非零项。
\end{theorem}
\begin{proofsketch}
	在 \(\chi\) 的表示中选择稳定满格，并按 \(C\) 的块分解其限制。关键的局部分解结论说，不属于 \(\operatorname{Br}_{\langle u\rangle}^G(b)\) 的那些块分量的每个不可分解项，都相对某个不包含 \(\langle u\rangle\) 的子群 \(Q\le C\) 投射。因为 \(u\) 在 \(C\) 中心，也不属于任何 \(C\)-共轭的 \(Q\)。Green 的相对投射特征标消失定理于是使这些项在 \(uv\)、\(v\in C_{p'}\) 处的特征标值为零。对整个诱导模，消失可从无固定陪集看出；对其不可分解直和项仍消失，则是这里真正需要的附加定理，不能由迹的可加性直接推出。这样，\(\chi(uv)\) 完全由 Brauer 像之下的其余块分量贡献；各块的普通特征标约化只产生本块的 Brauer 特征标，再用线性独立性就得到断言。

	这里省去的是满格的局部分解及上述直和项的特征标消失证明，二者均在特征零中进行，不能以 \(k\) 值矩阵迹为零代替。Nagao 的原文用相对迹理想与本原幂等元建立局部分解，在 Section 3 的 Theorem 1 得到所需的特征标等式，再在 Section 4 的 Theorem 2 推出广义分解数的限制；见\cite[Sections 3--4, Theorems 1--2]{Nagao1963GeneralizedDecompositionNumbers}。
\end{proofsketch}

不能将第二主定理读成“非零 Brauer 像所允许的广义分解数都非零”，也不能把 \(\langle u\rangle\) 自动当作原块的缺陷群。它只限制一张广义分解矩阵中哪些块外的条目必须为零。

\term[Brauer-disan-zhudingli]{Brauer 第三主定理}[Brauer's third main theorem]涉及主块的对应。下面陈述并证明其 Sylow 正规化子形式。
\begin{theorem}[Brauer 第三主定理的 Sylow 正规化子形式]\label{b5:thm:brauer-third-main-sylow}
	设 \(G\) 为有限群，\(k\) 为特征 \(p>0\) 的域，并且对 \(G\) 及其全部子群都是分裂域。令 \(P\) 为 Sylow \(p\)-子群，\(N=N_G(P)\)。记 \(b_0(G)\)、\(b_0(N)\) 为包含平凡模的主块幂等元，则 \(b_0(G)\) 的缺陷群为 \(P\)，且
	\[
		\operatorname{Br}_P^G(b_0(G))=b_0(N).
	\]
\end{theorem}
\begin{proof}
	若 \(D\) 是主块的缺陷群，由\cref{b5:cor:block-modules-relative-defect}，平凡模 \(k\) 相对 \(D\) 投射。在 \(\operatorname{End}_k(k)=k\) 上共轭作用平凡，所以 Higman 判据在此就是要求乘数 \([G:D]\) 非零。因 \(D\) 是 \(p\)-群，这等价于 \(D\) 为 Sylow \(p\)-子群。于是第一主定理保证 \(c=\operatorname{Br}_P^G(b_0(G))\) 是 \(kN\) 的块幂等元。

	增广 \(\varepsilon:kG\to k\) 给出平凡模的作用。块幂等元在平凡模上只能为零或一，且恰有主块取一。另一方面，\(P\) 的每个非单点共轭轨道的大小在 \(k\) 中为零，所以对 \(z\in(kG)^P\) 有 \(\varepsilon(\operatorname{Br}_P^G(z))=\varepsilon(z)\)。因此 \(\varepsilon(c)=\varepsilon(b_0(G))=1\)，只能有 \(c=b_0(N)\)。
\end{proof}
这个版本比较主块与 Sylow 正规化子的主块。一般第三主定理还描述任意 \(p\)-子群的中心化子中的主块及其 Brauer 对应；这里没有把这项更强的结论作为上述短证明的推论。

\subsection{Green 对应}
\label{b5:subsec:0806:03}

\term[Green-duiying]{Green 对应}[Green correspondence]选择的是不可分解模。它利用的识别条件并非简单性或维数，而是顶点。为避免混淆共轭范围，下列陈述中“顶点为 \(Q\)”表示可以选取 \(Q\) 本身为顶点；在 \(kH\) 模中，顶点的唯一性只到 \(H\)-共轭。
\begin{theorem}[Green 对应]\label{b5:thm:green-correspondence}
	设 \(G\) 为有限群，\(k\) 为特征 \(p>0\) 的域，\(Q\le G\) 为 \(p\)-子群，\(N_G(Q)\le H\le G\)。所有模均为有限维左模。
	\begin{enumerate}
		\item 若不可分解 \(kG\)-模 \(M\) 的顶点为 \(Q\)，则
		\[
			\operatorname{Res}_H^G M\cong f(M)\oplus X,
		\]
		其中恰有一个不可分解直和项 \(f(M)\) 的顶点为 \(Q\)，重数为一。\(X\) 的每个不可分解直和项都相对某个 \(H\cap{}^xQ\)、\(x\notin H\)，投射。
		\item 若不可分解 \(kH\)-模 \(V\) 的顶点为 \(Q\)，则
		\[
			\operatorname{Ind}_H^G V\cong g(V)\oplus Y,
		\]
		其中恰有一个不可分解直和项 \(g(V)\) 的顶点为 \(Q\)，重数为一。\(Y\) 的每个不可分解直和项都相对某个 \(Q\cap{}^xQ\)、\(x\notin H\)，投射。
		\item 这两个选择在同构类上互逆：\(g(f(M))\cong M\)、\(f(g(V))\cong V\)。
	\end{enumerate}
\end{theorem}
\symidx{\(f(M),g(V)\)：固定子群与顶点后的 Green 对应}
\begin{proof}
	若 \(x\notin H\)，则 \(Q\cap{}^xQ<Q\)，因为相等会迫使 \(x\in N_G(Q)\subseteq H\)。另一个交集 \(H\cap{}^xQ\) 却可能与 \(Q\) 同阶；真正需要的是它不与 \(Q\) 在 \(H\) 中共轭。若 \(H\cap{}^xQ={}^hQ\)、\(h\in H\)，两边同阶便有 \({}^xQ={}^hQ\)，从而 \(h^{-1}x\in N_G(Q)\subseteq H\)，仍是矛盾。

	先从顶点为 \(Q\) 的 \(kH\)-模 \(V\) 出发，取其源 \(S\)，写 \(\operatorname{Ind}_Q^H S=V\oplus W\)。由 Mackey 分解，\(\operatorname{Res}_H^G\operatorname{Ind}_Q^G S\) 中，代表元落在 \(H\) 的一项就是 \(\operatorname{Ind}_Q^H S=V\oplus W\)，其余项都相对 \(H\cap{}^xQ\)、\(x\notin H\)，投射。另一方面，\(\operatorname{Res}_H^G\operatorname{Ind}_H^G V\) 具有由单位双陪集给出的直和项 \(V\)，对 \(W\) 也如此。把这两式相加，并在 Krull--Schmidt 分解中消去 \(V\oplus W\)，便知
	\[
		\operatorname{Res}_H^G\operatorname{Ind}_H^G V\cong V\oplus V',
	\]
	且 \(V'\) 的全部不可分解项都相对上述 \(H\cap{}^xQ\) 投射。它们不可能有顶点 \(Q\)，故这里的 \(V\) 是唯一的该顶点项。

	分解 \(\operatorname{Ind}_H^G V\)，至少有一个不可分解项 \(M\) 的限制含 \(V\)。它相对 \(Q\) 投射；若它的顶点能选在 \(Q\) 的真子群中，限制后的 Mackey 分解会使 \(V\) 也相对一个阶更小的子群投射，矛盾。因此 \(M\) 的顶点是 \(Q\)。对另一项 \(M'\)，其限制的每个不可分解项都在 \(V'\) 中。取 \(M'\) 的顶点 \(Q'\le Q\)，并在限制到 \(H\) 后选取进一步限制到 \(Q'\) 含其源的一项 \(L\)。记 \(L\) 的一个顶点为 \(R\)。由 \(M'\) 相对 \(Q'\) 投射，再对其限制使用 Mackey 分解，得到 \(|R|\le|Q'|\)。另一方面，源是 \(\operatorname{Res}_{Q'}^H L\) 的直和因子；把 \(L\) 的相对 \(R\) 投射性代入 Mackey 分解，便使该源相对某个 \(Q'\cap{}^hR\)、\(h\in H\)，投射。源的顶点为 \(Q'\)，所以 \(Q'\le{}^hR\)；结合阶数不等式，得到 \(Q'={}^hR\)。因此可以选 \(Q'\) 本身为 \(L\) 的顶点。由于 \(L\) 是 \(V'\) 中的一项，某个 \(h\in H\) 满足
	\[
		{}^hQ'\le H\cap{}^yQ\quad(y\notin H),\qquad
		Q'\le Q\cap{}^{h^{-1}y}Q.
	\]
	后者是真子群。这证明诱导方向中的存在、唯一性及余项范围。

	反过来，设 \(M\) 是顶点为 \(Q\) 的不可分解 \(kG\)-模，取源 \(S\)。从 \(M\mid\operatorname{Ind}_Q^G S\) 及诱导传递性，在 \(\operatorname{Ind}_Q^H S\) 中选一项 \(V\)，使 \(M\mid\operatorname{Ind}_H^G V\)。\(V\) 只能有顶点 \(Q\)，否则 \(M\) 的顶点也会变小。由已证部分，\(M=g(V)\)。其限制到 \(H\) 后有一项进一步限制到 \(Q\) 含源 \(S\)，故该项具有顶点 \(Q\)；它只能是前面唯一的 \(V\)。其余项具有第一问指定的相对投射性。两方向的唯一性同时给出互逆关系。
\end{proof}

在限制方向中，余项的顶点不能一概说成“比 \(Q\) 小”：有些可能与 \(Q\) 在 \(G\) 中共轭，却不在 \(H\) 中共轭。诱导方向才由 \(Q\cap{}^xQ<Q\) 得到严格变小。这个区别也说明，不能只比较顶点的阶来挑出限制后的对应项。以上证明采用\cite[Theorem 11.6.4]{Webb2016FiniteGroupRepresentations}的源与 Mackey 分解方法。

\ProseParagraphBreak

Brauer 对应匹配块，Green 对应匹配指定顶点的不可分解模；两者都不是“限制整个模范畴就得到等价”的断言。特别地，诱导和限制一般仍带有额外直和项，简单模也可能对应到不简单的模。下面的计算直接追踪这些额外项。

\subsection{小群的块与对应计算}
\label{b5:subsec:0806:04}

先取特征二的 \(k\) 与 \(G=S_3\)，沿用\cref{b5:exm:s3-two-blocks}的 \(r=(123)\)、\(s=(12)\)、\(e=1+r+r^2\)、\(f=r+r^2\)。对 \(P=\langle s\rangle\)，有 \(N_G(P)=C_G(P)=P\)。已算出的
\[
	\operatorname{Br}_P^G(e)=1,\qquad \operatorname{Br}_P^G(f)=0
\]
现在给出实际的 Brauer 对应：缺陷为 \(P\) 的主块 \(kGe\) 对应 \(kP\) 的唯一块。矩阵块 \(kGf\) 的缺陷为一，不能把它也配给 \(kP\)；对它应取 \(N_G(1)=G\)，此时对应为恒等。

\ProseParagraphBreak

\(kP\cong k[t]/(t^2)\) 只有两个不可分解模 \(k\)、\(kP\)，其中只有 \(k\) 的顶点为 \(P\)。另一方面，\(kGe\cong k[t]/(t^2)\)，而矩阵块只有投射的简单模 \(S_2\)；所以 \(kG\) 中也只有平凡模的顶点为 \(P\)。Green 对应在这里把两边的平凡模配对，但诱导仍不等于这个对应项：
\[
	\operatorname{Ind}_P^{S_3}k\cong k\oplus S_2.
\]
验证可在三个陪集上的置换模 \(k^3\) 中进行。常向量构成平凡子模，坐标和为零的平面构成补模，因为 \(3=1\ne0\)；这个平面正是前面构造的 \(S_2\)。额外项 \(S_2\) 投射，顶点为一，符合 \(P\cap{}^xP=1\)、\(x\notin P\) 的结论。

\begin{example}\label{b5:exm:s3-three-block-local}
	设 \(k\) 为特征三的域，\(G=S_3=\langle r,s\mid r^3=s^2=1,\ srs=r^{-1}\rangle\)，\(P=\langle r\rangle\)。令 \(S_+\)、\(S_-\) 为平凡模和符号模，\(t=r-1\)、\(e_\pm=(1\pm s)/2\)。则 \(kG\) 只有主块，其缺陷群为 \(P\)，且
	\[
		J(kG)=(t),\quad J(kG)^3=0,\qquad P_\pm=kGe_\pm
	\]
	分别是 \(S_\pm\) 的投射覆盖。\(P_+\) 从顶到基座的根层为 \(S_+,S_-,S_+\)，\(P_-\) 的相应根层为 \(S_-,S_+,S_-\)。因此
	\[
		C=\begin{pmatrix}2&1\\1&2\end{pmatrix}.
	\]
\end{example}
\begin{proof}
	因为 \(t^3=r^3-1=0\)，且 \(sts=r^{-1}-1=-t+t^2\)，所以由 \(t\) 生成的双侧理想是 \(tkP\oplus tkP s\)，它的三次幂为零。商代数为 \(k(G/P)\cong kC_2\cong k\times k\)，是半单的；故该幂零理想恰为 \(J(kG)\)，商的两个简单模正是 \(S_+,S_-\)。

	\(e_+,e_-\) 是正交幂等元，和为一，所以 \(P_\pm\) 投射，且各以 \(e_\pm,te_\pm,t^2e_\pm\) 为基。\(J^jP_\pm/J^{j+1}P_\pm\) 由 \(t^je_\pm\) 的类张成；\(r=1+t\) 在它上面作用为一，而
	\[
		s t^je_\pm=(-t+t^2)^j s e_\pm
		\equiv \pm(-1)^j t^je_\pm\pmod{J^{j+1}P_\pm}.
	\]
	于是得到三层交替的符号。它们的顶均简单，故由\cref{b5:prop:projective-cover-top}，\(P_\pm\) 是相应投射覆盖且不可分解。一个不可分解模只能属于一个块，且它的全部合成因子都在同一块中。\(P_+\) 同时含两个简单模，所以这两个简单模属于同一块；每个非零块都含简单模，故全代数只有一个块，幂等元为一。对 \(P\) 有 \(\operatorname{Br}_P^G(1)=1\)，且 \(P\) 已为 Sylow 子群，因而它就是缺陷群。根层中的次数给出 Cartan 矩阵。
\end{proof}

在这个例子中 \(P\) 正规，故 \(N_G(P)=G\)。以 \(P\) 为顶点或缺陷群的对应都留在 \(G\) 本身，不能仅因“对应到了正规化子”就期待群阶下降。同一个群在特征二和特征三中的块数、缺陷群及对应计算均不同；决定这些区别的是特征与子群的位置，而不是普通特征标表的行数。

% * 状态：初稿已完成；尚未正式审核。
